% Source: https://artofproblemsolving.com/community/c3h1726527_spooky_logic_halloween_edition
% Stable identity: aops-topic-1726527-q01
Suppose I go trick-or-treating. Then:
\begin{itemize}
\item I get the most candy (statement 2), so (d) is \emph{false}.
\item Since inviting my best friend would mean not getting the most candy (statement 4), I do not invite my best friend: (c) is \emph{false}.
\item Not inviting my best friend means I invite my sister (statement 1). If I revealed a funny costume she would not come (statement 3), so I do not reveal that I am wearing a funny costume.
\end{itemize}
Nothing else follows. Whether I go to a haunted house is not determined: if I do, my friend gets scared, dares me to wear a serious costume, and I wear my funny costume. That is consistent with everything above, because \emph{wearing} the funny costume is not the same as \emph{revealing} it. If I don't go to the haunted house, all the statements also hold. So neither (a) nor (e) must be true. Finally, no statement ever produces a dare to wear a funny costume, so (b) need not be true.

Hence $\boxed{\text{none of (a)--(e) must be true}}$; in fact (c) and (d) are certainly false. Going trick-or-treating is consistent with all the statements: I invite my sister and not my best friend, get the most candy, and keep my costume a secret.
